\scanpage{85}
% archive foliation: 86

conditions supplémentaires étant sans doute équivalentes).\inserted{Le rang de $G$ comme groupe de BT s'appelle la \uline{hauteur} du groupe formel.}

c)\quad Groupes de BT ind-étales: par $G\quad T_p(G)$, \deleted{ces} la catégorie en question équivaut à celle des faisceaux $p$-adiques (étales) contants tordus \deleted{locale} sans torsion; si $S$ est connexe muni d'un point géométrique $\bar{s}$, cette catégorie équivaut à celle des représentations continues de $\pi_1(S,\bar{s})$ par des modules libres de type fini sur $\underline{Z}_p$.

Groupes de BT de t.m. ou ``\unsure{toroidaux|toroïdaux}'' : si $G(1)$, donc les $G(n)$, sont de type multiplicatif. Associant à un tel $G$ le système projectif des $G(n)^*$, on trouve \deleted{encore} une \inserted{anti}équivalence entre cette catégorie et la catégorie des faisceaux $p$-adiques constants tordus sans torsion. Composant avec le foncteur $F\mapsto F^{\inserted{\vee}}$, on trouve une équivalence au lieu d'une antiéquivalence, qui peut aussi s'expliciter comme \deleted{$G\mapsto\underline{\operatorname{Hom}}_{\mathrm{gr}}(\unsure{\mu_n},G)$}
\[
G\longmapsto\underline{\operatorname{Hom}}_{\mathrm{gr}}(\mu,G)
\overset{\mathrm{dfn}}{=}
\bigl(\underline{\operatorname{Hom}}_{\mathrm{gr}}(\mu(n),G(n))\bigr)_{n\geq1}.
\]
\margin{\unsure{k}}
\begin{center}(abus de notation)\end{center}
Noter que si les car. rés. sont $p$, alors les groupes de BT tor\deleted{\unsure{t}}oidaux sont \deleted{infinitésimaux} à fibres connexes (et réciproquement\inserted{s'il en est ainsi, car. rés sont $p$}).

d)\quad Groupes de BT comme extension d'un groupe ind-étale par un groupe de Lie formel.

Supposons d'abord $S$ \deleted{artinien (\unsure{ou}} réduit à un point\inserted{$s$}. Alors pour tout groupe fini $H$ sur $S$, $H^o$ est un sous-groupe, dépendant fonctoriellement de $H$. Si $H$ est fini localement libre, de même $H^o$, et $H/H^o$ est \deleted{\unsure{de}} un groupe fini localement libre (cf. SGA 3 $\mathrm{VI}_A$ \ldots), qui est étale. Ainsi $H$ s'écrit canoniquement et fonctoriellement comme extension d'un groupe étale par un groupe infinitésimal (tout fini localement libre). Si $k(s)$ est de car. $p\quad0$, alors procédant par dualité, on trouve un sous-groupe $G^{\mathrm{tm}}$ de $G^o=G^{\mathrm{inf}}$, qui est le plus grand sous-groupe de t.m de $G$, d'où une filtration à trois crans (qu'on avait eu le tort de ne définir précédemment que sur un \uline{corps} \deleted{de} de base). Passant maintenant à la
